\documentclass[10pt,a4paper]{amsart}
\usepackage[margin=19mm]{geometry}
\usepackage{amsmath,amssymb,mathtools}
\usepackage[scr=rsfs,cal=euler]{mathalfa}
\usepackage{xfrac}
\usepackage[unicode,hidelinks,bookmarks=false]{hyperref}
\newcommand{\ie}{i.e.\ }
\newcommand{\cf}{cf.\ }
\newcommand{\resp}{respectively\ }
\newcommand{\Subsection}{\subsection}
\providecommand{\nobreakdash}{\nobreak\hbox{-}}
\setlength{\emergencystretch}{2em}
\multlinegap0pt
\usepackage{bm}
\usepackage{bbm}
\usepackage{dsfont}
\usepackage{xspace}
\usepackage{cancel}
\usepackage{mathtools}
\usepackage{amsthm}
\makeatletter
\@ifundefined{theo}{\newtheorem{theo}{Theo}}{}
\@ifundefined{lemm}{\newtheorem{lemm}[theo]{Lemma}}{}
\makeatother
\title[Stability properties of adapted tangent sheaves on K\"ahler-]{Stability properties of adapted tangent sheaves on K\"ahler--Einstein log Fano pairs}
\author{Louis Dailly}
\date{Source: arXiv:2601.19608v1 (CC BY 4.0)}
\begin{document}
\maketitle

\noindent\emph{Source and licence.} Stability properties of adapted tangent sheaves on K\"ahler--Einstein log Fano pairs. Version arXiv:2601.19608v1 (CC BY 4.0), distributed under the Creative Commons Attribution 4.0 International licence, as recorded in the arXiv licence metadata for this exact version and shown on the abstract page https://arxiv.org/abs/2601.19608v1. Full rights notice: licenses/arxiv-2601.19608-CC-BY-4.0.txt.\par\smallskip

\noindent\textit{Editorial scope.} Counted unit: the Lemm whose source label is \texttt{Reduction of semistability to the resolution} in the article (source0.tex lines 827--831, source label \texttt{Reduction of semistability to the resolution}), delivered together with the article's own complete original proof of it (source0.tex lines 833--862, one \texttt{proof} environment of 30 lines). No mathematics is written, repaired or re-derived: statement and proof are the article's own text line for line. Required context retained uncounted: none. Every definition, display and label of those lines is retained unchanged. Omitted from this package and not redistributed: 1--826, 832--832, 863--1032. Nothing inside a retained excerpt is omitted or altered: every retained body is delivered exactly as the article's lines are written, so no omission receipt of any kind accompanies this package. Citations: the retained excerpts carry no citation command at all, so no bibliography entry of the article is transcribed and no reference is redistributed. Reference adaptation: none was needed and none was made; every reference inside the delivered unit resolves inside it, so no unresolved reference or citation is produced. Printed slips kept literal: no printed word, symbol, hypothesis or number of the counted statement or of its proof was altered. Layout and class: the article's own class is \texttt{amsart} with packages including inputenc, fontenc, babel, xy, fancyhdr, array, amsfonts, bbm, amsmath, amssymb, graphicx, stmaryrd, moreverb, geometry; the wrapper is the lane's pinned \texttt{amsart} editorial wrapper at 10pt on A4, which restates the theorem-like environments the retained text needs and adds the article's own macro definitions that the retained text needs (none), with extra wrapper packages (bm, bbm, dsfont, xspace, cancel, mathtools). The delivered numbering is the unit's own. Assets: no retained span contains a figure environment, a graphics-inclusion command, a PDF-page inclusion, a table or a TikZ picture; no image file is copied into this package, the package declares no asset dependency and no image file is redistributed with it. Licence: version arXiv:2601.19608v1 of this article is distributed under the Creative Commons Attribution 4.0 International licence, as recorded in the arXiv licence metadata for this exact version and shown on the abstract page https://arxiv.org/abs/2601.19608v1. A full rights notice is delivered at licenses/arxiv-2601.19608-CC-BY-4.0.txt. Characters: every retained excerpt is pure ASCII, so the delivered engine, XeLaTeX with its default Latin Modern text font, reports no missing character.
\begin{lemm}\label{Reduction of semistability to the resolution}
    With notation as above:\\
    $(i)$ if $\mathcal{T}_{\widehat{X}, \widehat{\Delta}, \widehat{f}}$ is $\widehat{f}^*\pi^*c_1(X, \Delta)$-semistable, then $\mathcal{T}_{X, \Delta, f}$ is $f^*c_1(X, \Delta)$-semistable,\\
    $(ii)$ if $\widehat{\mathcal{E}}_{\widehat{X}, \widehat{\Delta}, \widehat{f}}$ is $\widehat{f}^*\pi^*c_1(X, \Delta)$-semistable, then $\mathcal{E}_{X, \Delta, f}$ is $f^*c_1(X, \Delta)$-semistable.
\end{lemm}

\begin{proof}

We first prove $(i)$. By pulling back the orbifold tangent bundle above the orbifold locus of $(X, \Delta)$ through $\pi \circ \widehat{f} = f \circ \nu $, we have the following isomorphisms:
$$
(\nu^{*}\mathcal{T}_{X, \Delta, f})_{|\widehat{Y} \backslash E'} \simeq (\mathcal{T}_{\widehat{X}, \widehat{\Delta}, \widehat{f}})_{|\widehat{Y} \backslash E'}.
$$

We consider $\mathcal{F} \subseteq \mathcal{T}_{X, \Delta, f}$ a proper subsheaf. Let us denote $\widehat{\mathcal{F}}$ the saturation of $(\nu^*\mathcal{F})_{|\widehat{Y} \backslash E'}$ in $\mathcal{T}_{\widehat{X}, \widehat{\Delta}, \widehat{f}}$. Thanks to the projection formula, we have the following equalities of slopes:
$$
\mu_{\nu^*f^*c_1(X, \Delta)}(\widehat{\mathcal{F}}) = \mu_{f^*c_1(X, \Delta)}(\mathcal{F}),
$$
$$
\mu_{\nu^*f^*c_1(X, \Delta)}(\mathcal{T}_{\widehat{X}, \widehat{\Delta}, \widehat{f}}) = \mu_{f^*c_1(X, \Delta)}(\mathcal{T}_{X, \Delta, f}).
$$
From these equalities we deduce the semistability of $\mathcal{T}_{X, \Delta, f}$.

We now prove $(ii)$. We use the same arguments as in $(i)$. In order to do this, it suffices to observe that the projection formula applied to $\widehat{\mathcal{E}}_{\widehat{X}, \widehat{\Delta}, \widehat{f}}$ computes the slope of $\mathcal{E}_{X, \Delta, f}$. We observe that we have an equality of classes in $H^1(\widehat{Y} \backslash E', \Omega^{[1]}_{\widehat{X}, \widehat{\Delta}, \widehat{f}})$
$$
\nu^*f^*(-(K_X + \Delta)) = \widehat{f}^*\pi^*(-(K_X + \Delta)),
$$
we deduce:
$$
\left(\nu^{*}\mathcal{E}_{X, \Delta, f}\right)_{|\widehat{Y} \backslash E'} \simeq (\widehat{\mathcal{E}}_{\widehat{X}, \widehat{\Delta}, \widehat{f}})_{|\widehat{Y} \backslash E'}.
$$
This enables us to compute slopes by projection formula:
$$
\mu_{\nu^*f^*c_1(X, \Delta)}(\widehat{\mathcal{E}}_{\widehat{X}, \widehat{\Delta}, \widehat{f}}) = \mu_{f^*c_1(X, \Delta)}(\mathcal{E}_{X, \Delta, f}).
$$
We conclude the semistability of $\mathcal{E}_{X, \Delta, f}$ similarly to $(i)$.
\end{proof}

\end{document}
